\documentclass[11pt, xcolor={dvipsnames}, aspectratio=169]{beamer}

\usetheme{metropolis}
\usepackage{appendixnumberbeamer}
\usepackage{amsmath}
\usepackage{amssymb}
\usepackage{array}
\usepackage{mathtools}
\usepackage{caption}
\usepackage{fancyvrb}
\usepackage[separate-uncertainty]{siunitx}
\usepackage{overpic}
\usepackage{tikz}
\usepackage{booktabs}
\usepackage{multirow}
\usepackage[absolute, overlay]{textpos}
\usetikzlibrary{positioning}

\graphicspath{{logos/}{figures/}}

\titlegraphic{
  \includegraphics[height=1.1cm]{atlas_logo}
  \hfill
  \includegraphics[height=1.1cm]{uni_bonn_logo}
}

\setbeamercolor{background canvas}{bg=white}
\setbeamersize{text margin left=1.5em, text margin right=1.5em}

\setlength{\leftmargini}{4mm}
\setlength{\leftmarginii}{4mm}

% My definitions
\usepackage{defs}

\author{Christopher Deutsch}
\institute{University of Bonn}
\date{June XX, 2023}

\title{\vspace{1.5em}Search for Higgs Boson Pair Production in the \bbtautau
  Final State with the ATLAS Experiment at the LHC}

\begin{document}

\maketitle

% Aim for 25ish slides?
\begin{frame}[noframenumbering]{Notes: Content}
  30 minutes (talk) + 15 minutes (discussion):
  \begin{itemize}

  \item Conclusion \& Outlook (1 slide)

  \item Event Displays (2 slides)

  \end{itemize}
\end{frame}

% ------------------------------------------------------------------------------

\begin{frame}{Introduction}
  Why do particle physics?

  What is particle physics?

  What is the origin of elementary particle masses?
\end{frame}

% ------------------------------------------------------------------------------

\begin{frame}{The Standard Model}

  \begin{columns}
    \column{0.7\textwidth}
    \centering

    \vspace*{1em}

    \includegraphics[width=\textwidth, trim=0.9cm 0.3cm 0.2cm 1.5cm,
    clip]{theory/sm}

    \column{0.3\textwidth}

    \begin{itemize}
      \setlength{\itemsep}{1.5em}
    \item Built on principle of local gauge invariance of the Lagrangian
    \item How to incorporate masses of elementary particles?
    \end{itemize}
  \end{columns}

  % Built on the principle of local gauge symmetries of the Lagrangian

  % Gravitation not included

  % Problem: Masses of elementary particles

\end{frame}

% ------------------------------------------------------------------------------

\begin{frame}{Electroweak Symmetry Breaking}
  % Brout, Englert, Higgs -> EWSB

  \begin{columns}
    \column{0.65\textwidth}
    \centering

    % Add a (doublet) of complex scalar field with a special potential

    \begin{align*}
      V(\phi) = \mu^2 \phi^\dag \phi + \lambda (\phi^\dag \phi)^2 \\
      \mu^2 < 0
    \end{align*}

    \column{0.35\textwidth}
    \centering

    \includegraphics[width=\textwidth]{theory/potential}

    $\phi(x) = \phi_1(x) + i \, \phi_2(x)$
  \end{columns}

  \begin{align*}
    V(H) = \lambda v^2 H^2 + \lambda v H^3 + \frac{\lambda}{4} H^4
  \end{align*}
\end{frame}

% ------------------------------------------------------------------------------

\begin{frame}{Higgs Boson Discovery}

  \begin{columns}
    \column{0.5\textwidth}
    \centering

    % CERN Seminar: 4 July 2012
    \includegraphics[width=0.9\textwidth]{CERN_seminar}

    \column{0.5\textwidth}
    \centering

    \includegraphics[width=0.55\textwidth]{higgs_discovery/fig_002}

    \includegraphics[width=0.7\textwidth]{higgs_discovery/figaux_004a}
  \end{columns}
\end{frame}

% ------------------------------------------------------------------------------

\begin{frame}{The Higgs Potential}

  \begin{center}
    \begin{tabular}{ccccccc}
      $V(H)$ & $=$ & ${\color{Red} \lambda v^2 H^2}$ & $+$ & ${\color{Blue} \lambda v H^3}$ & $+$ & ${\color{Green} \frac{\lambda}{4} H^4}$ \\[1.2em]
             && \raisebox{1.85em}{\includegraphics[scale=0.8]{feynman_graphs_talk/higgs_massterm}}
             && \includegraphics[scale=0.8]{feynman_graphs_talk/higgs_hhh}
             && \includegraphics[scale=0.8]{feynman_graphs_talk/higgs_hhhh}
    \end{tabular}
  \end{center}
\end{frame}

% ------------------------------------------------------------------------------

\begin{frame}{Shortcomings of the Standard Model}
  Needed?
\end{frame}

% ------------------------------------------------------------------------------

\begin{frame}{The ATLAS Experiment at the LHC}
  \includegraphics[width=0.6\textwidth]{lhc/cern_complex_trim_otherexp_removed}%
  \includegraphics[width=0.6\textwidth]{atlas/atlas_overview}
\end{frame}

% ------------------------------------------------------------------------------

\begin{frame}{Non-Resonant \HH Production}
  \begin{columns}
    \column{0.5\textwidth}
    \centering

    \includegraphics[width=0.7\textwidth]{feynman_graphs/di_higgs_box}

    \column{0.5\textwidth}
    \centering

    \includegraphics[width=0.7\textwidth]{feynman_graphs/di_higgs_triangle}
  \end{columns}


  $\sigma_{\HH}^{gg\text{F}} = \SI{31.05}{\femto\barn}$

  $\sigma_{\HH}^{\text{VBF}} = \SI{1.7}{\femto\barn}$
\end{frame}

% ------------------------------------------------------------------------------

\begin{frame}{Resonant \HH Production}
  \includegraphics[width=0.35\textwidth]{feynman_graphs/di_higgs_resonant}
\end{frame}

% ------------------------------------------------------------------------------

\begin{frame}{Searches for \HH Production in \bbtautau Final States}
  Branching ratios of HH

  Hadhad / lephad

  Motivate the importance of hadronic tau lepton decays
\end{frame}

% ------------------------------------------------------------------------------

\begin{frame}{Tau Identification}
  What is tau identification?
\end{frame}

% ------------------------------------------------------------------------------

\begin{frame}{Tau Identification with RNN}
  How does it work with RNN?
\end{frame}

% ------------------------------------------------------------------------------

\begin{frame}{Performance of RNN-Based Tau Identification}
  It's great.

  \includegraphics[width=0.45\textwidth]{tauid/roc_incl_witherrors}
\end{frame}

% ------------------------------------------------------------------------------

\begin{frame}{Object and Event Selection 1}
\end{frame}

% ------------------------------------------------------------------------------

\begin{frame}{Object and Event Selection 2}
\end{frame}

% ------------------------------------------------------------------------------

\begin{frame}{Background Estimation: Overview}
\end{frame}

% ------------------------------------------------------------------------------

\begin{frame}{Background Estimation: Z+HF}
\end{frame}

% ------------------------------------------------------------------------------

\begin{frame}{Background Estimation: ttbarfakes}
  Possibly add second slide
\end{frame}

% ------------------------------------------------------------------------------

\begin{frame}{Background Estimation: multi-jet}
  Possibly add second slide
\end{frame}

% ------------------------------------------------------------------------------

\begin{frame}{Multivariate Analysis}
\end{frame}

% ------------------------------------------------------------------------------

\begin{frame}{Signal Extraction: non-resonant}
  \begin{columns}
    \column{0.5\textwidth}

    \column{0.5\textwidth}
    \centering
    \includegraphics[width=0.9\textwidth]{mva/prefit/Region_BMin0_incJet1_distSMBDT_J2_Y2015_DLLOS_T2_SpcTauHH_L0_Prefitlog}
  \end{columns}
\end{frame}

% ------------------------------------------------------------------------------

\begin{frame}{Signal Extraction: resonant}
\end{frame}

% ------------------------------------------------------------------------------

\begin{frame}{Statistical Analysis I}
\end{frame}

% ------------------------------------------------------------------------------

\begin{frame}{Statistical Analysis II}
  Somehow bring in uncertainties
\end{frame}

% ------------------------------------------------------------------------------

\begin{frame}{Results I}
  Possibly have to merge this with the previous part
\end{frame}

% ------------------------------------------------------------------------------

\begin{frame}{Results II}
  \includegraphics[width=0.5\textwidth]{results_res/resonant_upper_limits}
\end{frame}

% ------------------------------------------------------------------------------

\begin{frame}{Reinterpretation I}
  \includegraphics[width=0.4\textwidth]{feynman_graphs/di_higgs_effective}

  \begin{columns}
    \column{0.5\textwidth}
    \centering
    \includegraphics[width=0.9\textwidth]{self_coupling/hh_xsec}

    \column{0.5\textwidth}
    \centering
    \includegraphics[width=0.9\textwidth]{self_coupling/hh_mhh_vs_klam}
  \end{columns}
\end{frame}

% ------------------------------------------------------------------------------

\begin{frame}{Reinterpretation II}
  \includegraphics[width=0.5\textwidth]{self_coupling/klam_scan_result}
\end{frame}

% ------------------------------------------------------------------------------

\begin{frame}{Conclusion \& Outlook}
\end{frame}

% ------------------------------------------------------------------------------

\begin{frame}[standout]
  Backup
\end{frame}

% ------------------------------------------------------------------------------

\begin{frame}{Event Display: \hadhad Event}
  \centering
  \includegraphics[width=0.75\textwidth]{event_displays/hadhad}
\end{frame}

% ------------------------------------------------------------------------------

\begin{frame}{Event Display: \lephad Event}
  \centering
  \includegraphics[width=0.75\textwidth]{event_displays/lephad}
\end{frame}

% ------------------------------------------------------------------------------

\end{document}

%%% Local Variables:
%%% mode: latex
%%% TeX-master: "talk.tex"
%%% End:
